% ============================================================
% CAPITOLO 6 -- REFACTORING E VALIDAZIONE INCROCIATA
% ============================================================
\section{Refactoring e validazione incrociata}

Il punto 5 della consegna richiede di generare, tramite un modello
linguistico, quattro varianti rifattorizzate (C1-C4) di ciascuna
classe, con contesto crescente: C1 senza contesto, C2 con la suite
manuale (BB), C3 con l'aggiunta dei gap di coverage, C4 con
l'aggiunta dei test di mutation reinforcement (MT), dove disponibili.
Per ciascuna variante tutte le famiglie di test vengono rilanciate,
misurando pass/fail, copertura e mutation score rispetto a C0.
L'obiettivo è verificare se un refactoring assistito da LLM introduce
regressioni, e quali tecniche di test siano effettivamente in grado
di rilevarle: un refactoring che compili e superi tutti i test
esistenti non è automaticamente un refactoring sicuro, se nessuna di
quelle suite esercita la porzione di codice realmente modificata.
Per questo motivo, oltre all'esito aggregato pass/fail, viene
riportato per ciascuna classe anche il comportamento delle singole
famiglie di test: è a questo livello di dettaglio, e non nel solo
dato complessivo, che emergono le differenze più interessanti tra le
tecniche.

\subsection{Refactoring di \texttt{Resolver}}

Sono state generate quattro varianti rifattorizzate di
\texttt{Resolver}. C1 (nessun contesto) rinomina
\texttt{ErrorAction.error} in \texttt{errorType}, sostituisce
\texttt{Promote.isValid} con una tabella statica e riscrive
\texttt{toString()}, \texttt{noReorder()} e
\texttt{firstMatchingBranch} con stream e tabelle statiche. C2 (+BB)
applica le stesse modifiche di C1 ad eccezione del rename del campo
\texttt{error}, che il modello evita non appena riceve in input la
suite manuale. C3 (+gap coverage) si limita al solo refactoring di
\texttt{Promote.isValid}, l'unica modifica in una zona ben coperta
dalla suite manuale. C4 (+MT) è identica a C3, poiché
\texttt{Resolver} non ha una suite MT dedicata (Test Strength già
100\% con la sola suite BB).

La \tabref{tab:resolver-refactor-passfail} riassume l'esito per
famiglia:

\begin{itemize}
\item \textbf{BB (32 test)}: FAIL compilazione su C1, per il
riferimento diretto al campo \texttt{error} in un test di verifica su
\texttt{ErrorAction}. 32/32 da C2 in poi, senza modifiche al codice
di test.
\item \textbf{Randoop (campione di 500 test)}: FAIL compilazione su
C1, stesso motivo. Su C2, 5 test su 500 falliscono l'esecuzione: è
qui che è stato scoperto un bug di ricorsione, introdotto
semplificando la chiamata \texttt{Resolver.resolve(...)} in
\texttt{resolve(...)} -- la chiamata non qualificata veniva risolta
al metodo statico \texttt{resolve} della sottoclasse
(\texttt{RecordAdjust}, \texttt{WriterUnion} o \texttt{ReaderUnion}
dichiarano ciascuna la propria versione con la stessa firma) anziché
al dispatcher esterno. Corretto ripristinando la qualificazione
esplicita; 500/500 su C2-C4.
\item \textbf{LLM (58 test)}: FAIL compilazione su C1; 58/58 da C2 in
poi.
\item \textbf{EvoSuite (78 test)}: PASS su tutte le varianti, incluso
C1 -- genera le proprie sequenze per riflessione sui soli metodi
pubblici e non referenzia mai un campo interno per nome, quindi il
rename non lo tocca.
\end{itemize}

Nessuna delle quattro famiglie rileva il refactoring di
\texttt{toString()}, \texttt{noReorder()} e
\texttt{firstMatchingBranch} su C2: nessuna suite costruisce uno
scenario con più campi obbligatori mancanti insieme, l'unico caso in
cui la modifica sarebbe osservabile dall'esterno. È un punto cieco
condiviso da tutte le tecniche, non solo dalla suite manuale.

La \tabref{tab:resolver-refactor-delta} riporta i delta rispetto a
C0. Copertura: 84.7\%/68.1\% su C1 (solo suite ES compilante, quindi
non comparabile 1:1 con le altre varianti), 93.1\%/83.8\% su C2,
92.2\%/82.8\% su C3-C4. Mutation score: N/A su C1 (nessuna suite
eleggibile per PIT resta compilante), 51\% invariato da C2 in poi.
Complessità (proxy LOC/punti di decisione, in assenza di
un'analisi SonarCloud diretta sulle varianti): -15\% su C1-C2, dove
il refactoring più aggressivo introduce tabelle statiche e stream al
posto di cicli e confronti a cascata; -6\% su C3-C4, più
conservative.

Va inoltre segnalato un limite tecnico già discusso in
Sezione~\ref{sec:mutation}: PIT ed EvoSuite sono strutturalmente
incompatibili (conflitto tra il meccanismo di \emph{HotSwap} di PIT e
l'instrumentazione a runtime di EvoSuite), per cui il mutation score
qui riportato esclude sistematicamente la suite EvoSuite.

\subsection{Refactoring di \texttt{Utf8}}

Il design dei refactor per \texttt{Utf8} si è basato su un rischio
reale individuato nella classe durante l'analisi del codice sorgente:
\texttt{equals()} e \texttt{hashCode()} implementano un percorso
ottimizzato per array esattamente dimensionati, che richiede una
guardia esplicita sulla lunghezza effettiva del buffer; inoltre, i
metodi \texttt{set(String)} e \texttt{set(Utf8)} sono esercitati
dalle famiglie LLM, Randoop ed EvoSuite, ma da nessun test delle
famiglie BB e MT -- un possibile punto debole in caso di refactoring
proprio su quei metodi.

C1 (nessun contesto) rinomina \texttt{set} in
\texttt{setFromString}/\texttt{setFromUtf8}, unifica \texttt{equals()}
in un'unica chiamata a \texttt{Arrays.equals} e rimuove per errore la
guardia \texttt{bytes.length == length} da \texttt{hashCode()},
introducendo un bug latente sui buffer con capacità residua. C2 (+BB)
annulla il rename di \texttt{set}, ma il bug di \texttt{hashCode()}
resta: la suite BB non esercita mai lo scenario di buffer con
capacità residua. C3 (+gap coverage, che segnala esplicitamente
\texttt{setByteLength}, \texttt{equals} e \texttt{hashCode} come aree
a rischio) diventa conservativa, ripristina \texttt{equals()} e
\texttt{hashCode()} identici all'originale e sposta il refactoring su
\texttt{compareSequences()}. C4 (+MT) reintroduce con fiducia il
merge di \texttt{equals()} -- i test di boundary della suite MT sui
confini di lunghezza 7/8 lo confermano equivalente -- e rende
esplicita, anziché rimuovere, la guardia di \texttt{hashCode()}
tramite un metodo estratto \texttt{isExactlySized(...)}.

La \tabref{tab:utf8-refactor-passfail} riassume l'esito per famiglia:

\begin{itemize}
\item \textbf{BB (35 test)}: 35/35 su tutte le varianti. La suite non
referenzia mai \texttt{set} se non tramite l'interfaccia pubblica, e
il rename viene comunque annullato dal modello già in C2.
\item \textbf{MT (14 test)}: 13/14 su C1-C2 -- il fallimento è
esattamente il test che verifica il bug di \texttt{hashCode()};
14/14 su C3-C4, dove la guardia viene ripristinata o resa esplicita.
\item \textbf{Randoop (campione di 500 test)}: FAIL su C1 (rename di
\texttt{set}); 500/500 da C2. Il campionamento di input primitivi non
costruisce mai, nemmeno per caso, un'istanza tramite il costruttore
\texttt{package-private} \texttt{Utf8(String, int)} necessario a
esporre il bug di \texttt{equals()}.
\item \textbf{LLM (40 test)}: FAIL su C1; 40/40 da C2, stesso limite
di Randoop -- costruisce le proprie istanze solo tramite costruttori
pubblici.
\item \textbf{EvoSuite (63 test)}: FAIL su C1; 61/63 su C2 -- le 2
failure sono la prima evidenza del bug di \texttt{equals()}, raggiunto
costruendo per reflection un'istanza con lunghezza negativa tramite
\texttt{Utf8(String, int)}, accessibile a EvoSuite perché il test
generato si trova nello stesso package. 63/63 su C3 (refactoring
conservativo); 62/63 su C4, poiché il merge di \texttt{equals()}
viene reintrodotto.
\end{itemize}

I due difetti reali introdotti dal refactoring sono quindi
complementari e di difficoltà molto diversa: il bug di
\texttt{hashCode()} richiede solo API pubbliche, alla portata di
qualunque tecnica scelga lo scenario giusto (in pratica solo MT lo
fa); il bug di \texttt{equals()} richiede l'accesso a un costruttore
a visibilità ristretta, alla portata solo di una tecnica basata su
reflection strutturale come EvoSuite. Nessuna delle due tecniche
sarebbe stata sufficiente da sola.

La \tabref{tab:utf8-refactor-delta} riporta i delta rispetto a C0.
Copertura: 68.7\%/47.2\% su C1 (solo BB+MT compilanti), 100\%/97.2\%
su C2, 100\%/97.7\% su C3, 100\%/97.5\% su C4. Mutation score: N/A su
C1-C2, poiché la suite MT non è interamente verde -- precondizione
richiesta da PIT; 86\% su C3, invariato rispetto a C0; 85\% su C4,
lieve calo compatibile con il punto di decisione aggiuntivo introdotto
da \texttt{isExactlySized(...)}. Complessità: -11\% punti di decisione
su C1, la riduzione più marcata di tutte le varianti, ma anche
l'unica con un bug reale non rilevato dalla maggioranza delle
famiglie -- una controprova diretta del compromesso tra aggressività
del refactoring e sicurezza comportamentale.

\paragraph{Leggibilità dei test}
Un'osservazione qualitativa, verificata direttamente sul codice
sorgente generato: i test BB, MT e LLM hanno nomi di metodo
descrittivi e un commento esplicito di intento per ogni caso (es.
\texttt{testSetByteLengthExpansion}, con \texttt{@DisplayName} per la
suite LLM); Randoop ed EvoSuite usano invece naming puramente
sequenziale (\texttt{test0001}, \texttt{test00}) e nessun commento
sul comportamento verificato -- gli unici commenti presenti in
Randoop sono generati automaticamente dallo strumento per segnalare
un'eccezione attesa, non per documentare l'intento. Questa
differenza, verificata restando costante su tutte le varianti (la
stessa suite fisica viene ripuntata sul codice variante di turno,
anziché rigenerata da zero come nel report di riferimento della
stessa consegna), ha una ricaduta pratica diretta: i due interventi
di correzione descritti in Sezione~\ref{sec:mutation} sono stati
diagnosticabili solo perché il test coinvolto era leggibile; sarebbe
stato impraticabile farlo su una suite con la leggibilità di Randoop
o EvoSuite, dove l'intento originale del test non è mai esplicitato.

\subsection{Sintesi}

I risultati ottenuti sulle due classi convergono verso la stessa
conclusione, pur mostrando due facce diverse dello stesso problema.
Su \texttt{Resolver}, nessuna delle quattro tecniche rileva il
refactoring di \texttt{toString()}: è uno stesso punto cieco
condiviso da tutte, dovuto all'assenza di scenari con più campi
obbligatori mancanti insieme. Su \texttt{Utf8}, invece, il
refactoring introduce due difetti reali e distinti, e ciascuno risulta
visibile a una sola famiglia di test su cinque: il mutation testing
individua un valore di calcolo errato raggiungibile con le sole API
pubbliche, mentre EvoSuite individua una differenza di comportamento
raggiungibile solo sfruttando uno stato interno costruibile
esclusivamente tramite reflection o un costruttore riservato.

Questi risultati confermano empiricamente che nessuna singola tecnica
di test è sufficiente da sola a proteggere da regressioni introdotte
da un refactoring assistito da un modello linguistico con contesto
parziale. Ciascuna tecnica esplora lo spazio degli input secondo un
principio proprio -- rappresentatività funzionale per i test manuali,
sensibilità a mutazioni specifiche per il mutation testing,
casualità guidata dal feedback per Randoop, plausibilità sintattica
per gli LLM, massimizzazione della copertura strutturale per
EvoSuite -- e nessuno di questi principi, preso isolatamente,
garantisce la copertura degli scenari esplorati dagli altri. La
combinazione di più tecniche complementari, unita a un'analisi
mirata di copertura e mutation score per guidare eventuali interventi
correttivi, risulta pertanto necessaria per validare con affidabilità
un refactoring del codice.

Questo capitolo chiude idealmente il percorso avviato al Capitolo 3:
la suite manuale, la più debole in termini di copertura strutturale
grezza, si è rivelata qui l'unica in grado di intercettare un difetto
(\texttt{hashCode()}) proprio grazie ai casi limite individuati con
la tecnica Category-Partition, mentre le tecniche automatiche più
esplorative (EvoSuite) hanno raggiunto scenari -- come l'accesso a un
costruttore riservato -- che nessun progettista di test, umano o
guidato da un prompt, avrebbe verosimilmente pensato di includere. Il
valore aggiunto di un processo di verifica software non risiede
quindi in una singola tecnica "migliore delle altre", ma nella loro
combinazione mirata e nella capacità di interpretare criticamente i
dati di copertura e mutazione che ciascuna produce.